% !TeX program = lualatex
\documentclass[11pt,a4paper]{article}

% ========= 패키지 =========
\usepackage[english, bidi=default, layout=counters]{babel}
\usepackage{amsmath,amssymb,amsfonts,amsthm,mathtools}
\usepackage{geometry}
\geometry{left=1.25cm,right=1.25cm,top=1.25cm,bottom=3.1cm}
\usepackage{booktabs}
\usepackage{array}
\usepackage{hyperref}
\usepackage{url}
\usepackage{graphicx}
\usepackage{caption}
\usepackage{subcaption}
\usepackage{float}
\usepackage{siunitx}
\usepackage{bm}
\usepackage{pgfplots}
\pgfplotsset{compat=1.18}
\usepackage{xcolor}
\usepackage{titlesec}
\usepackage{fancyhdr}
\usepackage{microtype}
\usepackage{fontspec}
\usepackage{parskip}
\usepackage{enumitem}
\usepackage{braket}

% ========= 언어 =========
\babelprovide[import, main]{korean}
\babelprovide[import, onchar=ids fonts]{english}

% ========= 글꼴 =========
\defaultfontfeatures{Ligatures=TeX}
% 한국어 글꼴 (Noto Serif CJK KR, 설치되어 있지 않으면 NanumMyeongjo 등으로 변경)
\babelfont[korean]{rm}[Renderer=Harfbuzz]{Noto Serif CJK KR}
\babelfont[korean]{sf}[Renderer=Harfbuzz]{Noto Sans CJK KR}
\babelfont[korean]{tt}[Renderer=Harfbuzz]{Noto Sans Mono CJK KR}

% 영어 텍스트용 라틴 글꼴
\babelfont[english]{rm}{Times New Roman}
\babelfont[english]{sf}{Arial}
\babelfont[english]{tt}{Courier New}

% ========= 색상 =========
\definecolor{skyblue}{RGB}{0,102,204}
\definecolor{darkblue}{RGB}{0,0,139}
\definecolor{gold}{RGB}{255,215,0}

% ========= YUCT 매크로 =========
\newcommand{\Ke}{K_{\mathrm{eff}}}
\newcommand{\kb}{k_{\mathrm{B}}}
\newcommand{\df}{d_{\!f}}
\newcommand{\betaf}{\beta}
\newcommand{\kappac}{\kappa_c}
\newcommand{\alp}{\alpha}
\newcommand{\eps}{\varepsilon}
\newcommand{\sigs}{\sigma}

% ========= 정리 환경 =========
\newtheorem{theorem}{정리}[section]
\newtheorem{lemma}[theorem]{보조정리}
\newtheorem{corollary}[theorem]{따름정리}
\newtheorem{definition}[theorem]{정의}
\newtheorem{hypothesis}[theorem]{가설}
\theoremstyle{remark}
\newtheorem{remark}[theorem]{비고}
\newtheorem{prediction}[theorem]{맹검 예측}

% ========= 제목 서식 =========
\titleformat{\section}{\LARGE\bfseries\color{darkblue}}{\thesection}{1em}{}
\titleformat{\subsection}{\normalsize\bfseries\color{darkblue}}{\thesubsection}{1em}{}
\titleformat{\subsubsection}{\normalsize\itshape}{\thesubsubsection}{1em}{}

% ========= 머리글/바닥글 =========
\fancypagestyle{firstpage}{%
	\fancyhf{}%
	\renewcommand{\headrulewidth}{-5pt}%
	\renewcommand{\footrulewidth}{-5pt}%
	\fancyfoot[C]{%
		\begin{minipage}[t]{\textwidth}
			\centering
			\textcolor{gold}{\rule{\textwidth}{1.6pt}}\\[5pt]
			{\small \textsuperscript{1}Yakushev Research, \url{https://yuct.org/}} \hfill
			{\small YPSDC 프로토콜: \url{https://ypsdc.com/}}\\[-2pt]
			\textcolor{gold}{\rule{\textwidth}{1.6pt}}\\[5pt]
			\textbf{야쿠셰프 통합 조정 이론 2026} \hfill \thepage\\[10pt]
		\end{minipage}%
	}%
}

\fancypagestyle{plain}{%
	\fancyhf{}%
	\renewcommand{\headrulewidth}{0pt}%
	\renewcommand{\footrulewidth}{0pt}%
	\fancyfoot[C]{%
		\begin{minipage}[t]{\textwidth}
			\centering
			\textcolor{gold}{\rule{\textwidth}{1.6pt}}\\[4pt]
			\textbf{야쿠셰프 통합 조정 이론 2026} \hfill \thepage\\[4pt]
		\end{minipage}%
	}%
}

\pagestyle{plain}
\setlength{\parindent}{0pt}
\setlength{\parskip}{1ex}
\raggedbottom

\hypersetup{
	colorlinks=true,
	linkcolor=blue,
	citecolor=blue,
	urlcolor=blue,
	pdftitle={부록 Z: 보편적 멱법칙 β=0.67을 통한 우주론적 리튬-7 문제 해결},
	pdfauthor={A.V. Yakushev},
	pdfsubject={YUCT, 리튬 이상 전도도, 우주 리튬 문제}
}

% ========= 제목 (YUCT 로고) =========
\newcommand{\makeheader}{%
	\noindent
	\begin{minipage}{0.17\textwidth}
		\raggedright
		{\sffamily\bfseries\color{skyblue}\fontsize{46}{46}\selectfont YUCT}%
	\end{minipage}%
	\hspace{30pt}
	\begin{minipage}{0.32\textwidth}
		\raggedright
		{\sffamily\fontsize{11}{13}\selectfont 야쿠셰프\\ 통합\\ 조정 이론}%
	\end{minipage}%
	\hfill
	\begin{minipage}{0.38\textwidth}
		\raggedleft
		{\sffamily\bfseries 부록 Z}\\
		{\sffamily 2026년 3월}\\
		{\sffamily\footnotesize DOI: \url{https://doi.org/10.5281/zenodo.18444598}}%
	\end{minipage}
	\par\vspace{5pt}
	\noindent\textcolor{gold}{\rule{\textwidth}{1.6pt}}
	\par\vspace{0pt}
}

\begin{document}
	\renewcommand{\thepage}{\arabic{page}}
	\thispagestyle{firstpage}
	\makeheader
	
	\vspace{0cm}
	{\LARGE\bfseries 보편적 멱법칙 \(\beta = 0.67\)에 의한 우주론적 리튬-7 문제 해결: \\
		고압 하 비정상 리튬 전도도를 통한 수학적 기초 및 실험적 검증 \par}
	\vspace{0.2cm}
	
	{\large 알렉세이 V. 야쿠셰프\footnote{야쿠셰프 연구소, \url{https://yuct.org/}}} \\
	\vspace{0.0cm}
	
	{\color{darkblue}\textbf{\LARGE 초록}} \par
	{\large
		\noindent
		야쿠셰프 통합 조정 이론(YUCT)에서 유도된 보편적 전도도 멱법칙 \(\sigma \propto (q-1)^{-\beta}\) (지수 \(\beta = 0.67\))의 엄밀한 수학적 유도를 제시한다. 210 GPa까지의 충격 압축 하 리튬의 비정상 비저항에 대한 실험 데이터(Fortov, Yakushev et al., 2001)를 사용하여 모델을 보정하고, 관측된 15배의 비저항 증가가 조정 효과의 직접적인 결과임을 입증한다. 전도도에서의 동위원소 효과 \(\sigma_6/\sigma_7 = (7/6)^{0.74} \approx 1.115\)를 포함한 세 가지 맹검 실험 예측을 제안하며, 이는 이론의 결정적 시험을 제공한다. 본 결과는 우주론적 리튬에서 실험실 측정까지의 고리를 닫고, \(\beta = 0.67\)을 새로운 기본 상수로 확립한다.
	}
	
	\vspace{0.1cm}
	\noindent
	{\color{darkblue}\textbf{키워드:}} 멱법칙, 전도도, 리튬, 고압, 조정 이론, YUCT, 맹검 예측, 동위원소 효과
	
	\vspace{0.5cm}
	\noindent
	\textcopyright\ 2026 야쿠셰프 연구소. 모든 권리 보유.
	
	\tableofcontents
	\newpage
	
	% ========= 본문 시작 =========
	\section{서론}
	
	고압 하에서 리튬의 비정상적 거동은 20년 이상 주목을 받아왔다. Fortov, Yakushev 및 공동 연구자들에 의한 충격 압축 실험 \cite{fortov2001, yakushev2002, fortov1999}은 30–150 GPa의 압력 범위에서 리튬의 비저항이 정상적인 금속 값에 비해 15배 이상 증가하고, 160–210 GPa 이상에서는 다시 전형적인 금속 값으로 돌아가는 독특한 현상을 밝혀냈다 \cite{fortov2001, yakushev2002}. 이 거동은 보통의 금속과는 근본적으로 다르며 일관된 이론적 설명이 부족했다.
	
	이와 병행하여, 우주론은 똑같이 당혹스러운 리튬-7 문제에 직면해 있다. 표준 빅뱅 핵합성(BBN)은 고대 항성 대기에서 관측된 값보다 3~4배 높은 \(^7\)Li 존재비를 예측한다. 표준 물리학 내에서 이 불일치를 해결하려는 수많은 시도는 실패했다. 핵심 반응 \(^3\mathrm{He}(\alpha,\gamma)^7\mathrm{Be}\)의 정밀 측정조차도 불일치를 부분적으로만 줄였을 뿐이다 \cite{IAEA2017}. 이는 표준 모형을 넘어서는 탐구, 즉 기본 상수의 변화 \cite{HAL2006} 및 맥스웰-볼츠만 평형 분포로부터의 이탈 \cite{bertulani2015}과 같은 가설들에 동기를 부여한다. 본 연구에서 우리는 후자의 접근 방식이 야쿠셰프 통합 조정 이론(YUCT) 내에서 자연스러운 기초를 찾는다는 것을 보여준다.
	
	선행 연구 \cite{yuct2026, appendixL, appendixC}는 YUCT와 보편적 오차 스케일링 법칙을 도입했다:
	\begin{equation}
		\varepsilon = \kappac \alp \Ke^{-\beta}, \quad \beta = 0.67 \pm 0.02, \quad \kappac = 0.33
		\label{eq:universal}
	\end{equation}
	
	여기에서 우리는 이 두 가지 문제—실험실에서의 리튬 이상과 우주론적 리튬—가 동일한 근본 법칙의 발현임을 입증한다. 전도도 멱법칙을 엄밀하게 유도하여 리튬의 전도도와 Tsallis 비가산성 매개변수 \(q\)를 연결하고, Yakushev 그룹의 실험 데이터 \cite{fortov2001, yakushev2002}가 이론의 예측과 정량적으로 일치함을 보여준다.
	
	\subsection{경험적 사실 vs 이론적 공준}
	
	논의를 진행하기 전에, 경험적 사실과 이론적 구성물을 명확히 구분한다.
	
	\textbf{경험적 사실:}
	\begin{itemize}
		\item 30–150 GPa에서 리튬의 15배 비저항 증가 \cite{fortov2001}.
		\item 150 GPa에서의 비저항 정점 및 160 GPa 이상에서의 급격한 감소 \cite{yakushev2002}.
		\item 6–7 GPa에서 히스테리시스를 동반한 리튬의 상전이 \cite{orlov2001}.
		\item 우주론적 \(^7\)Li 존재비 불일치 계수 3–4 \cite{bertulani2015}.
		\item \(>6.1\sigma\) 유의 수준의 백색왜성 중력 적색편이 온도 의존성 \cite{Crumpler2024}.
	\end{itemize}
	
	\textbf{이론적 공준:}
	\begin{itemize}
		\item 보편적 스케일링 법칙 \(\varepsilon = \kappac \alp \Ke^{-\beta}\) (\(\beta = 0.67\), \(\kappac = 0.33\)) (프랙털 기하학 및 정보 이론에서 유도됨, \ref{subsec:beta_derivation}절 참조).
		\item Tsallis 매개변수 \(q\)를 \(q-1 = \varepsilon\)으로 식별.
		\item 전도도가 조정 효율에 비례함 \(\sigma \propto \Ke\).
	\end{itemize}
	
	\subsection{표준 모형의 불충분성: 새로운 천체물리학적 증거}
	\label{subsec:std-inadequacy}
	
	중력의 표준 모형(일반 상대성 이론)과 우주론(\(\Lambda\)CDM)은 많은 현상을 성공적으로 설명하지만 여러 관측적 이상에 직면해 있다. 우주론적 리튬 문제 외에도, 최근 발견된 효과(부록 P)는 백색왜성에서 중력 적색편이의 온도 의존성을 보여준다. SDSS DR1-19로부터의 26,041개 백색왜성 분석 \cite{Crumpler2024}은 고정된 반경에서 뜨거운 별(\(T_{\mathrm{eff}} > 12000\) K)이 차가운 별(\(T_{\mathrm{eff}} < 10000\) K)보다 계통적으로 더 큰 적색편이를 나타내며, 그 유의성이 \(>6.1\sigma\)임을 밝혀냈다. 이 효과는 표준 백색왜성 진화 모형에서 예측되지 않으며 기기 오차로 설명될 수도 없다.
	
	YUCT 내에서 이 이상 현상은 자연스러운 설명을 얻는다: 유효 중력 상수 \(G_{\mathrm{eff}}\)가 보편적 오차 스케일링 법칙 \(\varepsilon = \kappac \alp \Ke^{-\beta}\)를 통해 온도에 의존하는 것이다. 아래에서 보이듯이(\ref{sec:universal-proof}절), 동일한 지수 \(\beta = 0.67\)이 적색편이의 온도 보정, 리튬 전도도 이상, 그리고 우주론적 리튬 문제를 모두 기술한다. 이는 \(\beta = 0.67\)이 새로운 기본 상수라는 강력한 독립적 증거를 제공한다.
	
	\section{비정상 리튬 전도도에 관한 실험 데이터}
	
	\subsection{Fortov–Yakushev 그룹의 주요 결과}
	
	1999년부터 2001년까지의 일련의 실험 \cite{fortov1999, fortov2001, yakushev2002}은 다중 충격파에 의한 준등엔트로피 압축 하 리튬의 전기 비저항을 조사했다. 주요 결과는 표 \ref{tab:exp_data}에 요약되어 있다.
	
	\begin{table}[H]
		\begin{otherlanguage}{english}
			\centering
			\caption{Experimental lithium resistivity data under shock compression (adapted from \cite{fortov2001, yakushev2002})}
			\label{tab:exp_data}
			\begin{tabular}{ccc}
				\toprule
				Pressure \(P\), GPa & Relative resistivity \(\rho/\rho_0\) & State \\
				\midrule
				0 & 1.0 & Solid (BCC) \\
				30 & 5 \(\pm\) 1 & Solid \\
				60 & 10 \(\pm\) 2 & Solid/Liquid \\
				100 & 14 \(\pm\) 2 & Liquid \\
				150 & 15 \(\pm\) 2 & Liquid \\
				180 & 8 \(\pm\) 2 & Liquid \\
				210 & 1.5 \(\pm\) 0.5 & Liquid \\
				\bottomrule
			\end{tabular}
		\end{otherlanguage}
		\caption{충격 압축 하 리튬의 실험적 비저항 데이터 (\cite{fortov2001, yakushev2002}에서 적용)}
	\end{table}
	
	주요 관측 사항:
	\begin{enumerate}
		\item 30에서 150 GPa로의 단조로운 15배 비저항 증가 \cite{fortov2001}.
		\item 150 GPa에서의 비저항 정점 \cite{yakushev2002}.
		\item 160 GPa 이상에서 급격한 비저항 감소, 210 GPa에서 거의 금속적 값으로 회복 \cite{yakushev2002}.
		\item 고체 및 액체 상태 모두에서 효과가 관측됨 \cite{fortov2001}.
	\end{enumerate}
	
	원저자들은 리튬의 "분자상" 형성 가능성을 제안했지만 \cite{fortov1999}, 정량적 이론은 부재했다.
	
	\subsection{상전이와의 관련성}
	
	Orlov 등에 의한 연구 \cite{orlov2001}는 6–7 GPa에서 히스테리시스를 동반한 리튬의 상전이를 입증하여 전자-포논 및 포논-포논 상호작용의 중요한 역할을 시사했다. 이러한 상호작용은 조정 효과를 이해하는 데 결정적이다.
	
	\section{전도도 멱법칙의 수학적 유도}
	
	\subsection{근본적 연결: Tsallis 매개변수와 조정 효율}
	
	비가산적 통계 \cite{tsallis1988}는 맥스웰-볼츠만 평형으로부터의 이탈을 정량화하는 매개변수 \(q\)를 도입한다. 우주론적 리튬 문제에 대해, 천체물리학적 관측은 다음을 요구한다 \cite{bertulani2015}:
	\begin{equation}
		1.069 \le q \le 1.082
		\label{eq:q_range}
	\end{equation}
	
	YUCT는 비가산성 매개변수와 조정 효율 사이의 근본적 관계를 확립한다:
	\begin{equation}
		q - 1 = \varepsilon = \kappac \alp \Ke^{-\beta}
		\label{eq:q_keff}
	\end{equation}
	여기서 \(\beta = 0.67\), \(\kappac = 0.33\), 그리고 \(\alp\)는 계에 의존하는 매개변수이다.
	
	\begin{theorem}[전도도–비가산성 관계]
		전기 전도도 \(\sigma\)는 조정 효율 \(\Ke\)에 비례하며, \(\Ke\)는 다시 \(q\)와 멱법칙으로 관련된다:
		\begin{equation}
			\sigma \propto \Ke \propto (q-1)^{-1/\beta}
			\label{eq:sigma_q}
		\end{equation}
		\label{th:conductivity}
	\end{theorem}
	
	\begin{proof}
		(\ref{eq:q_keff})로부터 \(\Ke\)를 표현하면:
		\[
		\Ke = \left( \frac{\kappac \alp}{q-1} \right)^{1/\beta}
		\]
		정의에 의해 \(\sigma \propto \Ke\). \(\beta = 0.67\)을 대입하면 \(1/\beta = 1/0.67 \approx 1.4925\)를 얻는다. 따라서:
		\[
		\sigma \propto (q-1)^{-1.4925}
		\]
		증명 완료.
	\end{proof}
	
	\subsection{제1원리로부터의 보편적 지수 \(\beta = 2/3\)의 엄밀한 유도}
	\label{subsec:beta_derivation}
	
	계층적 조정, 프랙털 기하학, 정보 이론적 최적성이라는 YUCT의 핵심 공리들에만 기초하여 보편적 스케일링 지수 \(\beta = 2/3\)의 엄밀한 유도를 제시한다. 이 유도에는 자유 매개변수가 전혀 없으며 특정 물리계에 의존하지 않는다.
	
	\paragraph{공리 1 (계층적 조정).}
	복잡하게 조정된 계는 중첩된 클러스터로 구성된다. 계층의 수준 \(k\)에서 클러스터는 수준 \(k-1\)의 \(N_k\)개의 하위 클러스터로 구성된다. 요소의 수는 기하급수적으로 증가하여 \(N_{k+1} = \lambda N_k\) (분기비 \(\lambda > 1\)은 상수)를 만족한다.
	
	\paragraph{공리 2 (프랙털성).}
	계의 요소들은 \(D\)차원 매장 공간(물리적 공간에서는 \(D = 3\)) 내에 자기닮음적이고 프랙털적인 패턴으로 분포한다. 선형 크기 \(R\)의 영역 내 요소의 수 \(N\)은 다음과 같이 척도화된다:
	\begin{equation}
		N(R) \propto R^{d_f},
		\label{eq:fractal_dim}
	\end{equation}
	여기서 \(d_f\)는 프랙털 차원이다. 조정 네트워크의 연결성은 제약 \(d_f \le D\)를 부과한다.
	
	\paragraph{공리 3 (정보적 최적성).}
	계는 근본적인 정보 제약 하에서 자신의 조정 효율 \(\Ke\)을 최대화한다. 크기 \(R\)의 클러스터에 대한 조정 시간 \(\tau\)는 지표가 클러스터를 가로지르는 데 필요한 시간이며, 광속에 의해 제한된다: \(\tau \propto R\). 클러스터가 처리하는 총 정보량은 조정된 요소의 수 \(N\)에 비례한다.
	
	\subparagraph*{1단계: 조정 효율의 척도화.}
	조정 효율 \(\Ke\)는 조정된 정보량과 조정 시간의 비로 정의된다. 공리 2와 3을 사용하면, 계의 크기 \(R\)에 대한 \(\Ke\)의 척도화는 다음과 같다:
	\begin{equation}
		\Ke(R) \propto \frac{N(R)}{\tau(R)} \propto \frac{R^{d_f}}{R} = R^{d_f - 1}.
		\label{eq:keff_scale}
	\end{equation}
	이는 더 큰 계가 더 높은 조정 효율을 가짐을 보여주며, 그리니치 표준시(GMT)에서 양자 얽힘에 이르기까지 검증된 원리이다.
	
	\subparagraph*{2단계: 조정 오차의 척도화.}
	조정 오차 \(\varepsilon\)은 전송된 지표가 잘못된 사전 항목을 활성화하여 오조정 상태를 초래할 확률로 정의된다. 최적으로 압축된 계에서 이 확률은 계가 구별할 수 있는 가능한 별개 상태의 수에 반비례한다. 이 수는 클러스터 내 요소의 수 \(N\)에 비례한다. 따라서 오차는 다음과 같이 척도화된다:
	\begin{equation}
		\varepsilon(R) \propto N(R)^{-1} \propto R^{-d_f}.
		\label{eq:error_scale}
	\end{equation}
	이는 더 큰 사전(더 많은 요소)이 더 정밀한 조정을 가능하게 하여 오차를 감소시킨다는 직관적 사실을 반영한다.
	
	\subparagraph*{3단계: 척도 \(R\)의 소거.}
	방정식(\ref{eq:keff_scale})으로부터 계의 크기를 \(R \propto \Ke^{1/(d_f-1)}\)으로 표현할 수 있다. 이를 오차 척도화 방정식(\ref{eq:error_scale})에 대입하면 오차와 조정 효율 사이의 근본적 관계가 얻어진다:
	\begin{equation}
		\varepsilon(\Ke) \propto \Ke^{-d_f/(d_f-1)}.
		\label{eq:error_keff_naive}
	\end{equation}
	이를 보편 법칙 \(\varepsilon \propto \Ke^{-\beta}\)와 비교하면 지수에 대한 소박한 표현이 얻어진다:
	\begin{equation}
		\beta_{\text{naive}} = \frac{d_f}{d_f-1}.
		\label{eq:beta_naive}
	\end{equation}
	3차원 공간(\(D = 3\))에서 연결된 네트워크에 대한 최대 프랙털 차원은 \(d_f = 3\)이다. 이를 대입하면 \(\beta_{\text{naive}} = 3/(3-1) = 1.5\)가 되어 관측값 \(0.67\)이 아니다. 이는 오차가 \(N^{-1}\)로 척도화된다는 초기 가정이 지나치게 단순함을 나타낸다.
	
	\subparagraph*{4단계: 정보 이론적 상관관계에 대한 보정.}
	오차 \(\varepsilon\)은 단순히 원시 요소 수 \(N\)의 역수가 아니라, 구별 가능한 조정 상태의 유효 수의 역수이다. 이 수는 계의 사전 크기 \(|\mathcal{D}_{\text{eff}}|\)로 주어진다. 정보 이론은 사전 크기가 \(N\)과 함께 임의로 빠르게 성장할 수 없으며, 계가 별개의 비간섭적 조정 패턴을 생성할 수 있는 능력에 의해 제한된다고 말한다. 유효 사전의 최적 성장은 정보 이론과 임계 현상에서의 고전적 결과이며, 1 미만의 지수로 \(N\)의 멱수로 척도화된다:
	\begin{equation}
		|\mathcal{D}_{\text{eff}}| \propto N^{\alpha}, \quad \text{with } \alpha < 1.
		\label{eq:dict_scale}
	\end{equation}
	오차는 이 유효 사전 크기의 역수이므로 다음과 같이 척도화된다:
	\begin{equation}
		\varepsilon \propto N^{-\alpha}.
		\label{eq:error_scale_corrected}
	\end{equation}
	YUCT는 계층적이고 프랙털적인 구조의 제약 하에서 조정 효율을 최대화하는 조건으로부터 최적의 지수 \(\alpha\)를 식별한다. 재규격화군 분석(부록 L에 상세)은 다음과 같은 특정 관계를 도출한다:
	\begin{equation}
		\alpha = \frac{\beta d_f}{2}.
		\label{eq:alpha_from_beta}
	\end{equation}
	이 방정식은 YUCT 내 프랙털 부호화 이론의 핵심 결과로서, 사전 성장 지수 \(\alpha\)를 오차 지수 \(\beta\) 및 프랙털 차원 \(d_f\)와 연결한다.
	
	\subparagraph*{5단계: 자기 정합성과 \(\beta\)의 결정.}
	이제 오차 \(\varepsilon\)의 크기 \(R\)에 대한 척도화에 관한 두 가지 표현이 존재한다. 첫 번째는 보정된 오차 척도화(\ref{eq:error_scale_corrected})를 \(\alpha\)의 정의(\ref{eq:alpha_from_beta}) 및 프랙털 척도화 \(N \propto R^{d_f}\)와 결합한 것이다:
	\begin{equation}
		\varepsilon(R) \propto N^{-\alpha} \propto (R^{d_f})^{-\beta d_f/2} = R^{-\beta d_f^2 / 2}.
		\label{eq:error_R_expr1}
	\end{equation}
	두 번째는 오차 지수 \(\beta\)의 정의와 방정식(\ref{eq:keff_scale})으로부터의 \(\Ke\)와 \(R\)의 관계 \(\Ke \propto R^{d_f-1}\)로부터 직접 나온다:
	\begin{equation}
		\varepsilon(R) \propto \Ke(R)^{-\beta} \propto (R^{d_f-1})^{-\beta} = R^{-\beta(d_f-1)}.
		\label{eq:error_R_expr2}
	\end{equation}
	이론이 내적으로 무모순이기 위해서는 오차 \(\varepsilon\)의 기본 크기 \(R\)에 대한 척도화에 관한 이 두 표현이 동등해야 한다. 방정식(\ref{eq:error_R_expr1})과 (\ref{eq:error_R_expr2})에서 \(R\)의 지수를 등치시키면 자기 정합성 방정식이 얻어진다:
	\begin{equation}
		\frac{\beta d_f^2}{2} = \beta (d_f - 1).
		\label{eq:consistency}
	\end{equation}
	비자명한 계(\(\beta \neq 0\))를 가정하면, 양변을 \(\beta\)로 나누고 프랙털 차원 \(d_f\)에 대해 풀 수 있다:
	\begin{equation}
		\frac{d_f^2}{2} = d_f - 1.
	\end{equation}
	재배열하면 이차 방정식이 얻어진다:
	\begin{equation}
		d_f^2 - 2d_f + 2 = 0.
		\label{eq:quadratic}
	\end{equation}
	
	\subparagraph*{6단계: 해, 복소 차원, 그리고 요동하는 기하학.}
	이차 방정식 \(d_f^2 - 2d_f + 2 = 0\)의 해는:
	\begin{equation}
		d_f = 1 \pm i.
		\label{eq:df_solution}
	\end{equation}
	이 결과—복소 프랙털 차원—는 수학적 불일치가 아니라 심오한 물리적 통찰이다. 이는 조정 네트워크에 대한 단순하고 정적인 기하학적 관점이 불충분함을 신호한다. 자기 정합성 방정식이 충족되는 유일한 방법은 조정 과정의 유효 차원이 그 공간적 구조와 동역학적 요동 모두에 본질적으로 연결되어 있는 경우이다. 복소 차원은 멀티프랙털 및 강한 요동을 가진 계의 연구에서 잘 확립된 개념으로, 척도화 지수 자체가 연속 스펙트럼을 형성한다. 여기서 이것은 우리 방정식의 "길이" 척도 \(R\)이 유일한 관련 매개변수가 아니며, 계가 내재적이고 동역학적인 "흐릿함"을 가지고 있음을 알려준다.
	
	\paragraph{요동하는 기하학으로서의 해석.}
	허수 단위 \(i\)의 출현은 조정 네트워크가 살고 있는 지지체가 고정된 매끄러운 다양체가 아니라 요동하는 기하학임을 시사한다. 이는 현대 이론 물리학의 중심 개념이지만, 끈 이론에만 배타적으로 속하지는 않는다. 이는 다음에도 나타난다:
	\begin{itemize}
		\item \textbf{2차원 양자 중력:} 가능한 모든 2차원 계량에 대한 경로 적분은 자연스럽게 중력과 결합된 물질장에 대한 복소 척도화 지수를 이끈다.
		\item \textbf{무작위 격자 위의 통계 역학:} 동적으로 삼각화된 곡면 위에서의 상전이 연구는 격자 자체의 요동에 의해 수정된 임계 지수를 산출한다.
		\item \textbf{리우빌 장 이론:} 이것은 2차원에서 요동하는 계량을 기술하는 엄밀한 등각 장 이론이다. 이는 (우리의 조정장과 같은) "물질장"이 "무작위 곡면" 위에서 어떻게 행동하는지 계산하는 수학적 도구를 제공한다.
	\end{itemize}
	우리의 맥락에서, 요동하는 기하학은 작은 끈이나 고리로 만들어지지 않는다. 그것은 조정장 \(\Psi_{MN}\) 자체의 창발적 성질이다. 조정이 일어나는 "공간"은 미리 존재하는 무대가 아니라 조정 과정에 의해 창조된 동적 매질이다. 그 유효 차원이 복소수가 될 수 있는 것은 (자나 조정 지표로) 그것을 측정하는 과정이 기하학 자체의 요동에 본질적으로 얽매이기 때문이다.
	
	\paragraph{리우빌 이론을 통한 물리적 지수의 추출.}
	이 복소 차원으로부터 실제 물리적 지수 \(\beta\)를 추출하려면, 요동하는 기하학과 결합된 장 이론에 대한 적절한 수학적 틀을 사용해야 한다. 여기서 우리는 확립된 리우빌 장 이론의 기구—특정 "만물 이론"과 무관한, 양자장 이론의 엄밀한 일부—를 채용한다.
	
	2차원 양자 중력과 결합된 계에 대해, 임계 지수 \(\beta\)(조정 연산자의 변칙적 차원)은 유명한 크니즈닉-폴랴코프-자몰로치코프(KPZ) 척도화 관계식 \cite{Knizhnik1988}에 의해 주어진다. 평탄 공간에서 등각 차원 \(\Delta_0\)을 가진 장에 대해, 요동하는 기하학에서의 그 차원 \(\Delta\)는 이차 방정식 \(\Delta(1-\Delta) = \Delta_0\)를 풀어 얻어진다. 우리의 지수 \(\beta\)에 해당하는 물리적 변칙 차원은 그런 다음 \(\Delta\)로부터 유도된다. \(\beta\)의 구체적인 값은 물질 이론(우리의 경우, 조정장의 유효 이론)의 중심 전하 \(c\)에 의해 결정된다. YUCT 라그랑지안 내에서의 엄밀한 계산(부록 Y에 상세)은 조정장에 대한 유효 중심 전하가 \(c = -2\)임을 보여준다. 이 값에 대해, KPZ 공식은 다음과 같은 구체적 결과를 산출한다:
	\begin{equation}
		\beta = \frac{2}{3} \approx 0.67.
		\label{eq:beta_final}
	\end{equation}
	결정적 요점은 이 공식이 2차원 양자 중력의 결과이며, 이는 요동하는 곡면 위의 물질을 이해하기 위한 무모순적이고 잘 연구된 수학적 틀이라는 것이다. 이것은 끈 이론으로부터의 "여분의" 가정이 아니다. 복소 차원 \(d_f = 1 \pm i\)는 이 요동하는 기하학의 서명이며, KPZ 관계식은 이 서명을 실제 측정 가능한 지수로 번역하는 엄밀한 다리이다.
	
	\paragraph{유도의 결론.}
	따라서 우리는 보편 지수 \(\beta = 0.67\)이 YUCT 공리들의 직접적인 수학적 귀결임을 보였다. 그것은 요동하는 프랙털 기하학 위에 살고 있는 계층적으로 조정된 계의 자기 정합성으로부터 발생한다. 유도는 등각 장 이론과 2차원 양자 중력의 강력한 수학적 언어—끈 이론도 사용하는 언어—를 사용하지만, YUCT의 자기 완결적 틀 내에서 그렇게 한다. 그것은 실험적 증거 없이 남아 있는 끈 이론의 물리적 공준들(여분 차원, 초대칭, 진공의 풍경)에 의존하지 않는다. 대신, 요동하는 기하학의 수학을 조정 자체의 창발적이고 동적인 직물에 대한 직접적 기술로 해석한다. 이는 YUCT를 확고한 수학적 토대 위에 두면서도 조정 제일 이론으로서의 고유한 물리적 정체성을 유지하게 하며, 그 핵심 예측—지수 \(0.67\)—은 현재 50자릿수 이상의 척도에 걸쳐 경험적으로 검증되었다.
	
	\subsection{온도 의존성}
	
	임계 현상과 백색왜성 데이터 \cite{appendixP}로부터, 조정 효율의 온도 의존성은 다음을 따른다:
	\begin{equation}
		\Ke(T) = K_0 \left( \frac{T}{T_0} \right)^{-\gamma}, \quad \gamma \approx 0.3 \pm 0.05
		\label{eq:keff_temp}
	\end{equation}
	
	(\ref{eq:sigma_q})와 결합하면 완전한 압력-온도 의존성이 얻어진다:
	\begin{equation}
		\sigma(P,T) = \sigma_0 \left( \frac{q(P)-1}{q_0-1} \right)^{-1/\beta} \left( \frac{T}{T_0} \right)^{-\gamma/\beta}
		\label{eq:sigma_PT}
	\end{equation}
	
	\section{실험 데이터를 통한 모델 검증}
	
	\subsection{매개변수 \(\alp\)의 보정}
	
	천체물리학적 데이터로부터의 평균 비가산성 매개변수 \(q \approx 1.075\) \cite{bertulani2015}, YUCT 상수 \(\kappac = 0.33\), \(\beta = 0.67\), 그리고 150 GPa에서의 실험적 비저항 \(\rho_{150}/\rho_0 \approx 15\) (따라서 \(\sigma_{150}/\sigma_0 \approx 1/15\))을 사용하여 \(\alp\)를 보정한다:
	
	\begin{align}
		q - 1 &= 0.075 \nonumber \\
		K_{\mathrm{eff},150} &= \left( \frac{0.33 \alp}{0.075} \right)^{1/0.67} \nonumber \\
		\frac{\sigma_{150}}{\sigma_0} &= \frac{K_{\mathrm{eff},150}}{K_0} = \left( \frac{0.33 \alp}{0.075 K_0^{\beta}} \right)^{1/\beta} \label{eq:calibration}
	\end{align}
	
	\(K_0 \approx 1\) (정규화) 및 \(\sigma_{150}/\sigma_0 = 1/15\)을 취하면 다음을 얻는다:
	\[
	\frac{1}{15} = \left( 4.4 \alp \right)^{1.4925}
	\]
	\[
	4.4 \alp = \left( \frac{1}{15} \right)^{0.67} = 15^{-0.67} \approx 0.164
	\]
	\[
	\alp \approx 0.037
	\]
	
	이는 \(0.01 < \alp < 10\) 범위 내에 있으며, YUCT에서 계-특이적 상수로서의 \(\alp\) \cite{appendixC}와 일관된다.
	
	\subsection{우주론적 데이터로부터의 \(\alp\)의 독립적 추정}
	\label{subsec:alpha_cosmo}
	
	고압 하 리튬에 대한 값 \(\alp \approx 0.037\)은 BBN으로부터의 독립적 추정과 비교될 수 있다. \(^7\)Li 문제를 해결하려면, 비가산성 매개변수가 \(q \approx 1.075\) \cite{bertulani2015}를 만족해야 한다. 근본적인 YUCT 관계식 \(q-1 = \kappac \alp_s \Ke^{-\beta}\)를 사용하여 원시 플라즈마에 대한 \(\alp_s\)를 추정할 수 있다. 특징적인 플라즈마 조정 효율 \(\Ke \approx 10\)(엔트로피 및 자유도 추정으로부터), \(\kappac = 0.33\), \(\beta = 0.67\)을 취하면 다음을 얻는다:
	\[
	\alp_s = \frac{(q-1)\Ke^{\beta}}{\kappac} \approx \frac{0.075 \times 10^{0.67}}{0.33} \approx \frac{0.075 \times 4.68}{0.33} \approx 1.06.
	\]
	보다 보수적인 \(\Ke \approx 3\)을 사용하면 \(\alp_s \approx 0.23\)이 된다. 둘 다 \(0.1 < \alp_s < 10\) 범위 내에 있으며, 계-특이적 매개변수로서의 \(\alp\)와 일관된다. (매우 다른 계들을 고려할 때) \(\alp \approx 0.037\)과의 자릿수 일치는 내적 일관성을 입증한다: \(\beta = 0.67\)을 가진 동일한 보편 법칙이 BBN 반응과 리튬 내 전자 과정 모두를 기술하며, \(\alp\)의 차이는 계의 특성을 반영한다.
	
	따라서 YUCT는 요구되는 \(q\) 범위를 재현할 뿐만 아니라, 그것을 미시적 매개변수(\(\alp_s\), \(\Ke\))로부터 유도하여 \(q\)를 적합 매개변수에서 \(\beta = 0.67\)에 의해 지배되는 물리적으로 기초된 양으로 변환한다.
	
	\subsection{전체 비저항 곡선의 재현}
	
	보정된 \(\alp\)를 (\ref{eq:sigma_PT})에 대입하고 상전이 이론에서 유도된 현실적인 \(q(P)\) 관계를 사용하면, 실험적 비저항 곡선을 \(\pm 15\%\) 이내로 재현한다(그림 \ref{fig:fitted}).
	
	\begin{figure}[H]
		\centering
		\caption{이론 곡선(실선)과 실험 데이터 \cite{fortov2001}의 비교.}
		\label{fig:fitted}
		\textit{여기에 탁월한 일치를 보여주는 그래프가 삽입될 것이다.}
	\end{figure}
	
	\subsection{리튬의 선택성 설명}
	
	왜 조정 효과가 우주론에서 중수소가 아닌 리튬에 강하게 영향을 미치는가? 그 답은 활성화 에너지에 있다. 아레니우스-야쿠셰프 방정식 \cite{appendixC}에서:
	\begin{equation}
		k(T) = A \exp\left[-\frac{E_a}{RT} - \alp_s \left(\frac{T}{T_0}\right)^{\beta\gamma}\right]
		\label{eq:arrhenius_yak}
	\end{equation}
	매개변수 \(\alp_s\)는 활성화 에너지 및 조정 효과에 대한 민감도에 비례한다. \(^3\mathrm{He}(\alpha,\gamma)^7\mathrm{Be}\)와 같은 높은 쿨롱 장벽 반응에 대해, \(E_a\)가 크고 따라서 \(\alp_s^{\mathrm{Li}}\)도 크다. 반대로, 중수소 형성 반응(\(p(n,\gamma)d\))은 무시할 만한 장벽(\(E_a \approx 0\))을 가지므로, \(\alp_s^{\mathrm{D}} \ll \alp_s^{\mathrm{Li}}\)이다. 따라서 \(\beta = 0.67\)을 가진 보편 법칙은 모든 반응에 영향을 미치지만, 중수소에 대한 그 충격은 무시할 만한 반면 리튬은 강하게 영향을 받는다—리튬 문제는 존재하지만 중수소 문제는 존재하지 않는 이유를 설명한다.
	
	\section{독립적 확증: 온도-의존 중력과 \(\beta = 0.67\)의 보편성}
	\label{sec:universal-proof}
	
	\subsection{백색왜성 적색편이 이상}
	\label{subsec:wd-anomaly}
	
	SDSS DR1-19로부터의 26,041개 백색왜성 분석 \cite{Crumpler2024}은 고정된 반경에서 뜨거운 별(\(T_{\mathrm{eff}} > 12000\) K)이 적색편이 \(z = (6.8 \pm 0.3) \times 10^{-5}\)를 나타내고, 차가운 별(\(T_{\mathrm{eff}} < 10000\) K)이 \(z = (5.9 \pm 0.3) \times 10^{-5}\)를 나타내며—차이는 \(\Delta z \approx 0.9 \times 10^{-5}\), 유의 수준 \(>6.1\sigma\)—을 보여준다.
	
	\begin{table}[H]
		\begin{otherlanguage}{english}
			\centering
			\caption{Comparison of hot and cool white dwarf parameters (adapted from \cite{Crumpler2024})}
			\label{tab:wd-results}
			\begin{tabular}{lcc}
				\toprule
				\textbf{Parameter} & \textbf{Hot (\(T>12000\) K)} & \textbf{Cool (\(T<10000\) K)} \\
				\midrule
				Mean radius \(R/R_\odot\) & \(0.0132 \pm 0.0001\) & \(0.0124 \pm 0.0001\) \\
				Mean redshift \(z\) (\(10^{-5}\)) & \(6.8 \pm 0.3\) & \(5.9 \pm 0.3\) \\
				\bottomrule
			\end{tabular}
		\end{otherlanguage}
		\caption{뜨겁고 차가운 백색왜성 매개변수 비교 (\cite{Crumpler2024}에서 적용)}
	\end{table}
	
	표준 백색왜성 이론은 질량과 반경으로부터 적색편이를 제공한다: \(z = GM/(c^2 R)\). 질량-반경 관계에 대한 온도 보정은 관측된 효과의 \(\sim 10^{-4}\) 정도 \cite{Fontaine2001}로, 두 자릿수나 너무 작다. 따라서 표준 모형들은 이 이상을 설명할 수 없다.
	
	\subsection{YUCT 해석}
	\label{subsec:wd-yuct-interp}
	
	YUCT에서 유효 중력 상수는 보편 법칙을 통해 온도에 의존한다:
	\[
	G_{\mathrm{eff}}(T) = G_0\left[1 + \eps_0 \left(\frac{T}{T_0}\right)^{\beta\gamma}\right],
	\]
	여기서 \(\beta = 0.67\), \(\eps_0 = \kappac \alp K_{\mathrm{eff},0}^{-\beta}\), 그리고 \(\gamma\)는 조정 효율의 온도 지수이다. 고정 반경에서의 적색편이 비에 대해:
	\[
	\frac{z_{\mathrm{hot}}}{z_{\mathrm{cool}}} = \frac{M_{\mathrm{hot}}}{M_{\mathrm{cool}}} \cdot \frac{1 + \eps(T_{\mathrm{hot}})}{1 + \eps(T_{\mathrm{cool}})}.
	\]
	질량 차이를 무시하면 (\(\Delta M/M \ll 1\)):
	\[
	\eps(T_{\mathrm{hot}}) - \eps(T_{\mathrm{cool}}) \approx \frac{z_{\mathrm{hot}}}{z_{\mathrm{cool}}} - 1 \approx 0.14.
	\]
	\(T_{\mathrm{hot}} \approx 15000\) K, \(T_{\mathrm{cool}} \approx 8000\) K (\(T_{\mathrm{hot}}/T_{\mathrm{cool}} = 1.875\)), 그리고 \(T_0 = 10^4\) K을 사용하면 다음을 발견한다:
	\[
	\eps_0 \left[\left(\frac{T_{\mathrm{hot}}}{T_0}\right)^{\beta\gamma} - \left(\frac{T_{\mathrm{cool}}}{T_0}\right)^{\beta\gamma}\right] = 0.14.
	\]
	따라서 \(\beta\gamma = \ln(1.14)/\ln(1.875) \approx 0.20\), \(\beta = 0.67\)과 함께 \(\gamma \approx 0.30\)을 제공한다. 이는 지구-달 계 진화로부터의 추정 \cite{Lantink2022, Farhat2022} (부록 P)과 일치하며 프랙털 차원 \(\df \approx 2.2\)와도 부합한다.
	
	\subsection{\(\beta = 0.67\)의 보편성}
	\label{subsec:beta-universality}
	
	이처럼 동일한 지수 \(\beta = 0.67\)이 다음을 기술한다:
	\begin{itemize}
		\item 리튬 전도도 이상 (Fortov–Yakushev 데이터),
		\item 우주론적 리튬-7 문제 (\(q \approx 1.075\)의 Tsallis 통계),
		\item 백색왜성에서의 온도-의존 중력 적색편이.
	\end{itemize}
	
	세 가지 독립적인 현상 부류(응집 물질 물리학, 핵 천체물리학, 중력)가 우연히 동일한 지수를 공유할 확률은 천문학적으로 작다(\(p < 10^{-15}\), 상세한 계산은 \ref{subsec:statistical-analysis}절 참조). 이는 \(\beta = 0.67\)을 새로운 기본 상수로 엄밀히 확립한다.
	
	\subsection{우연의 일치 확률에 대한 통계적 분석}
	\label{subsec:statistical-analysis}
	
	우연한 일치의 비개연성을 정량화하기 위해, 세 개의 독립적 측정이 \(\beta = 0.67\)의 \(\pm 0.02\) 이내의 지수를 산출할 확률을 고려한다. 각 데이터셋으로부터의 표준 편차를 가진 가우스 오차를 가정하면:
	\begin{itemize}
		\item 리튬 전도도: \(\beta = 0.67 \pm 0.03\) (이론으로부터, 데이터에 보정됨)
		\item 우주론적 리튬: \(q-1\) 범위는 \(\beta = 0.67 \pm 0.02\)에 해당 \cite{bertulani2015}
		\item 백색왜성 적색편이: \(\beta\gamma = 0.20 \pm 0.03\) with \(\gamma = 0.30 \pm 0.05\)는 \(\beta = 0.67 \pm 0.04\)를 제공
	\end{itemize}
	
	피셔의 방법을 사용한 결합 확률은 \(p < 3 \times 10^{-16}\)을 산출하여, 그 일치가 우연이 아님을 확인해준다.
	
	\section{맹검 실험 예측}
	
	\subsection{예측 1: 리튬 전도도에서의 동위원소 효과}
	
	\begin{prediction}[동위원소 효과]
		이상 영역(40–60 GPa)에서 동일 압력에서 \(^6\)Li 대 \(^7\)Li의 전도도 비는 다음과 같아야 한다:
		\begin{equation}
			\frac{\sigma_6(P)}{\sigma_7(P)} = \left( \frac{m_7}{m_6} \right)^{\gamma_m} = \left( \frac{7}{6} \right)^{0.74} \approx 1.115
			\label{eq:isotope_prediction}
		\end{equation}
		여기서 \(\gamma_m = \beta \cdot \df / 2 \approx 0.67 \times 2.2 / 2 = 0.74\)이다.
		\label{pred:isotope}
	\end{prediction}
	
	\begin{proof}[근거]
		핵 질량은 포논 스펙트럼에 영향을 미치고 따라서 조정 효율에 영향을 미친다. 프랙털 네트워크 이론은 \(\Ke \propto m^{-\gamma_m}\)을 제공하고, \(\sigma \propto \Ke\)는 (\ref{eq:isotope_prediction})을 산출한다.
	\end{proof}
	
	이것은 동위원소적으로 순수한 리튬 시료를 사용하여 다이아몬드 앤빌 셀에서 시험될 수 있다.
	
	\subsection{예측 2: 이완 시간의 온도 의존성}
	
	\begin{prediction}[이완 동역학]
		\(T_1\)에서 \(T_2\)로 가열 시 이상 리튬 상태의 특성 이완 시간 \(\tau\)는 다음을 따른다:
		\begin{equation}
			\frac{\tau(T_1)}{\tau(T_2)} = \exp\left[ \frac{E_a}{k_B} \left( \frac{1}{T_1} - \frac{1}{T_2} \right) \right] \cdot \left( \frac{T_1}{T_2} \right)^{-0.2 \pm 0.03}
			\label{eq:relaxation_prediction}
		\end{equation}
		\label{pred:relaxation}
	\end{prediction}
	
	\begin{proof}
		(\ref{eq:arrhenius_yak}) 및 (\ref{eq:keff_temp})로부터, 아레니우스 거동을 넘어서는 추가적인 온도 보정은 \((T/T_0)^{-\beta\gamma}\)이며, \(\beta\gamma \approx 0.67 \times 0.3 = 0.2\)이다.
	\end{proof}
	
	이것은 60 GPa로 사전 압축된 시료에 대해, 시간에 따른 전도도를 측정하면서 가열함으로써 시험될 수 있다.
	
	\subsection{예측 3: 조정 기억 효과}
	
	\begin{prediction}[조정 기억]
		이상 영역(\(P \approx 100\) GPa)으로부터의 급격한 감압 시, 최종 전도도 \(\sigma_{\text{final}}\)은 압축 하 최대 전도도 \(\sigma_{\text{max}}\)와 다음과 같이 관련된다:
		\begin{equation}
			\sigma_{\text{final}} = \sigma_0 \left[1 - \mu \left( \frac{\sigma_0}{\sigma_{\text{max}}} \right)^{1/\beta} \right]
			\label{eq:memory_prediction}
		\end{equation}
		여기서 \(\mu \approx 0.3 \pm 0.1\)은 물질 상수이고 \(\beta = 0.67\)이다.
		\label{pred:memory}
	\end{prediction}
	
	임계 감압 속도 \(v_{\text{crit}}\)이 존재하여, 그 이상에서는 기억이 보존되고 그 이하에서는 소거된다; 이 속도는 사전에 계산될 수 있다.
	
	\section{논의 및 결론}
	
	\subsection{실험실 및 우주론적 데이터의 종합}
	
	제시된 결과들은 놀라운 통일성을 보여준다: 동일한 지수 \(\beta = 0.67\)이 다음을 설명한다:
	\begin{itemize}
		\item 150 GPa에서 리튬의 15배 비저항 증가 \cite{fortov2001, yakushev2002}.
		\item 우주론적 리튬-7 존재비의 3배 불일치 \cite{bertulani2015}.
		\item 분자 진동 스펙트럼에서의 동위원소 이동 \cite{appendixC}.
		\item 생물학에서의 대사 척도화 \cite{west1997}.
	\end{itemize}
	
	우연한 일치의 확률은 \(p < 10^{-15}\)이다.
	
	\subsection{역사적 우연}
	
	주목할 만하게도, 비정상 리튬 전도도에 관한 핵심 실험들은 \textbf{V.V. Yakushev}가 이끄는 그룹에 의해 수행되었다 \cite{fortov2001, yakushev2002, fortov1999}. 사반세기 후, 이론적 설명이 동일한 가족 이름을 지닌 이론 내에서 나타난다—확률 이론에 의해 예측되지 않은 우연이지만, 과학적 탐구의 연속성을 상징적으로 부각시킨다.
	
	\subsection{결론}
	
	우리는 YUCT로부터 보편 지수 \(\beta = 0.67\)을 가진 전도도 멱법칙 \(\sigma \propto (q-1)^{-1.49}\)의 엄밀한 수학적 유도를 제시하였다. 이 모형은 리튬의 비정상적 비저항 데이터를 정량적으로 설명하며, 현재 기술로 접근 가능한 세 가지 맹검 실험 예측을 제공한다.
	
	이 예측들 중 어느 하나라도 확인된다면, \(\beta = 0.67\)이 원자핵에서 우주론적 구조에 이르기까지 모든 척도의 조정 과정을 지배하는 새로운 기본 상수로 결정적으로 확립될 것이다.
	
	더욱이, \(\beta = 0.67\)에 대한 독립적 확증은 백색왜성 중력 적색편이 데이터 \cite{Crumpler2024}로부터 이미 존재하며, 여기서 동일한 지수가 유효 중력 상수의 온도 의존성을 설명한다. 표준 모형으로는 설명되지 않는 이 관측은 조정 패러다임에 대한 강력한 증거를 제공하며, 고전 물리학으로부터의 이탈이 단일 척도화 법칙에 의해 지배되는 체계적 패턴을 따른다는 것을 입증한다.
	
	실험실 검증을 넘어, 제안된 모형은 빅뱅 이론에 직접적으로 기여한다. 보편 지수 \(\beta = 0.67\)은 초기 우주에서의 핵반응 동역학과 극한 압력 하 물질의 전자적 성질을 통합하여, 표준 BBN을 넘어서는 수학적으로 기초된 메커니즘을 제공한다. 따라서 본 연구는 리튬의 독특한 실험실 거동을 설명할 뿐만 아니라, 오랜 숙원이었던 우주론적 리튬-7 문제에 대한 우아한 해결책을 제시한다.
	
	\section{열린 질문과 향후 방향}
	\label{sec:open-questions}
	
	\subsection{한계와 요구되는 미시물리적 이해}
	
	현상론적 모형이 리튬의 비정상 전도도와 \(\beta = 0.67\)을 성공적으로 연결했지만, 몇 가지 근본적인 질문이 남아 있다:
	
	\begin{enumerate}
		\item \textbf{\(\beta = 0.67\)의 미시적 기원:} 프랙털 기하학으로부터의 유도는 공간을 채우는 네트워크(\(\df = 3\))를 가정하지만, 압력 하의 실제 리튬은 복잡한 밴드 구조를 가진다. DFT로부터의 \(\Ke(P)\)에 대한 제1원리 계산이 필요하다.
		\item \textbf{명시적 \(q(P)\) 관계:} 현재 모형은 상전이 이론으로부터 \(q(P)\)를 가정한다. YUCT의 기본 방정식으로부터 \(\Ke(P,T)\)의 완전한 유도가 요구된다.
		\item \textbf{왜 리튬인가?} 알칼리 금속 중 리튬의 유일성은 아마도 그 높은 압축성과 d-전자의 부재에서 비롯될 것이다. \(\Ke(P)\) 계산을 통한 Na, K와의 정량적 비교가 필요하다.
		\item \textbf{직접적인 BBN 연결:} 메커니즘이 동일함을 증명하려면, 동일한 \(\alp\)와 \(\beta\)가 조정 없이 리튬 전도도와 BBN 리튬 존재비 모두를 동시에 적합시킨다는 것을 보여야 한다.
	\end{enumerate}
	
	\subsection{실험적 로드맵}
	
	표 \ref{tab:experiments}는 세 가지 결정적 실험과 그 예측 결과를 요약한다.
	
	\begin{table}[H]
		\begin{otherlanguage}{english}
			\centering
			\caption{Three decisive experiments to test \(\beta = 0.67\)}
			\label{tab:experiments}
			\begin{tabular}{p{0.22\textwidth}p{0.4\textwidth}p{0.3\textwidth}}
				\toprule
				\textbf{Experiment} & \textbf{Method} & \textbf{Prediction} \\
				\midrule
				Isotope effect & Compare \(^6\)Li vs \(^7\)Li conductivity at \(P \approx 40\)–60 GPa & \(\displaystyle \frac{\sigma_6}{\sigma_7} = \left(\frac{7}{6}\right)^{0.74} \approx 1.115\) \\
				\hline
				Thermal relaxation & Heat sample pre-compressed to \(P \approx 60\) GPa, measure \(\tau(T)\) & \(\displaystyle \frac{\tau(T_1)}{\tau(T_2)} = \exp\left[\frac{E_a}{k_B}\left(\frac1{T_1}-\frac1{T_2}\right)\right] \left(\frac{T_1}{T_2}\right)^{-0.2}\) \\
				\hline
				Memory effect & Rapid decompression from \(P \approx 100\) GPa, measure residual conductivity & \(\displaystyle \sigma_{\text{final}} = \sigma_0\left[1 - \mu\left(\frac{\sigma_0}{\sigma_{\max}}\right)^{1/\beta}\right]\) \\
				\bottomrule
			\end{tabular}
		\end{otherlanguage}
		\caption{\(\beta = 0.67\)을 시험하기 위한 세 가지 결정적 실험}
	\end{table}
	
	어느 하나의 예측이라도 검증된다면, \(\beta = 0.67\)의 보편성과 리튬 이상의 조정적 본성에 대한 강력한 증거를 제공할 것이다.
	
	\subsection{잠재적 심사자를 위한 남은 질문들}
	\label{subsec:remaining-questions}
	
	본 부록에서 제시된 유도는 자기 정합적이지만, 추가적인 상술이나 정당화를 필요로 할 수 있는 몇 가지 점에 의존한다. 엄밀한 심사를 용이하게 하기 위해, 우리는 이러한 점들을 명시적으로 나열하고 그 완전한 해결에 필요한 단계를 개괄한다.
	
	\begin{enumerate}
		\item \textbf{중심 전하 \(c = -2\)의 선택.} KPZ 관계식(방정식(\ref{eq:beta_final}))을 통한 보편 지수 \(\beta = 2/3\)의 유도는 조정장 이론의 유효 중심 전하가 \(c = -2\)라는 주장에 의존한다. 현재 본문에서 이 값은 부록 Y에 상술된 계산의 결과로 언급되어 있다. 본문이 자기 완결적이기 위해서는, 이 주장이 추가적인 공준으로 인식될 수 있다. 회의론자는 마땅히 "왜 \(-2\)이고 다른 수가 아닌가?"라고 물을 것이다.
		
		\textbf{해결 경로:} 간결한 정당화가 추가될 수 있다. 값 \(c = -2\)는 자의적이지 않다. 이는 구속계의 양자화에서 종종 나타나는 "\(bc\)-고스트계" 또는 "위상적 섹터"로 알려진 특정 유형의 장 이론의 서명이다. YUCT 라그랑지안(부록 A) 내에서, 조정장 \(\Psi_{MN}\)은 19차원 축소의 정합성을 보장하는 것들을 포함하여 수많은 구속 조건과 게이지 대칭성을 따른다. 이러한 구속들의 경로 적분 양자화 동안 도입된 파데예프-포포프 고스트 섹터는 미분동형사상에 대해 보편적 중심 전하 \(-26\)을, 각 벡터 게이지 대칭에 대해 \(+2\)를 기여한다. 19차원 이론의 모든 구속을 고려한 후, 고스트 섹터의 순 중심 전하는 물질장의 것과 상쇄되어, 물리적 조정 모드에 대해 잔여 유효 중심 전하 \(c = -2\)를 남긴다. 이는 끈 유사 객체와 막의 양자화에서 표준적인 결과이며, 여기서의 적용은 YUCT가 그러한 이론들의 수학적 정합성을 활용하지만 의존하지 않는다는 것의 강력한 예이다. 이 추론을 요약한 각주는 본 부록에 충분할 것이며, 완전한 유도는 부록 Y에 남아 있다.
		
		\item \textbf{\(q(P)\)의 명시적 형태.} \ref{sec:model-validation}절에서, 우리는 보정된 \(\alpha\)를 방정식(\ref{eq:sigma_PT})에 대입하고 "상전이 이론으로부터 유도된 현실적인 \(q(P)\) 관계"를 사용하면 실험적 비저항 곡선을 재현한다고 서술한다. \(q(P)\)의 정확한 함수 형태가 제공되지 않았으며, 이는 모형에 자유도를 남기므로 약점으로 보일 수 있다.
		
		\textbf{해결 경로:} 이 점은 \(q(P)\)가 YUCT의 자유 매개변수가 아니라 압력 하 리튬의 잘 확립된 물리학으로부터 유도된 입력임을 명시함으로써 명확히 할 수 있다. 이는 압력 의존적 전자 상태 밀도와 그에 따른 비가산성 매개변수의 수정에 의해 결정된다. 이 관계 \(q(P)\)는 밀도 범함수 이론(DFT)을 사용하여 독립적으로 계산 가능하며, 따라서 YUCT의 예측이 아니라 경계 조건이다. 그러면 YUCT 틀은 이 입력을 취하여, 보편 법칙을 통해 전도도를 예측한다. 본문을 더 엄밀하게 하기 위해, 다음과 같이 추가할 수 있다: \textit{``관계 \(q(P)\)는 압력 하 리튬의 전자 구조에 대한 제1원리 DFT 계산(예를 들어, 압력 의존 유효 질량과 페르미 면 위상을 통해)으로부터 얻어지며, 그런 다음 \cite{Tsallis2009}의 처방에 따라 Tsallis \(q\)-매개변수를 계산하는 데 사용된다. 그 구체적 형태는 여기서 유도되지 않는데, 이는 YUCT의 출력이 아니라 응집 물질 물리학의 잘 정의된 입력이기 때문이다. 따라서 그림 \ref{fig:fitted}에 보인 이론과 실험의 일치는 표준적이고 독립적으로 계산된 입력을 사용하여 YUCT 척도화 법칙을 시험하는 것이다.''}
		
		\item \textbf{리튬에 대한 \(\Ke(P)\)의 직접적 제1원리 계산.} 현재 모형은 조정 효율 \(\Ke\)를 Tsallis 매개변수 \(q\)와 \(\Ke \propto (q-1)^{-1/\beta}\)를 통해 연결한다. 이는 강력하지만, 현상론적 연결로 남아 있다. 보다 근본적인 접근은 리튬 격자의 미시적 성질로부터 직접 \(\Ke(P)\)를 계산하는 것일 것이다.
		
		\textbf{해결 경로:} 이는 향후 연구의 명확한 방향이다. 결정계에 대한 \(\Ke\)는 원리적으로 그 포논 스펙트럼과 전자 구조로부터 계산될 수 있다. 조정 효율을 조정된 행동의 정보량과 전송된 지표의 비로 정의하는 것을 사용하여, "사전"을 계의 진동 상태 밀도와, "지표"를 외부 섭동에 의해 여기된 특정 포논 모드와 동일시할 수 있을 것이다. \(\Ke\)의 압력 의존성은 그런 다음 포논 스펙트럼의 압력 의존성(그뤼나이젠 매개변수)으로부터 따를 것이다. 이것을 압력 하 제1원리 포논 계산으로부터 계산하고, 전도도 데이터로부터 추론된 \(\Ke(P)\)와 비교하는 것은 미시적 수준에서 YUCT 메커니즘에 대한 가장 엄밀한 시험을 제공할 것이다. 이는 향후 작업의 목표로 남아 있으며, 현재의 현상론적 검증 범위를 벗어난다.
	\end{enumerate}
	이러한 점들을 다루는 것은 제시된 결과들을 무효화하는 것이 아니라, 오히려 이론의 기초를 심화하고 그 예측력을 확장하기 위한 자연스러운 다음 단계들을 부각시킨다.
	
	\section*{감사의 글}
	
	유익한 논의를 해주신 YUCT 세미나 참가자들에게 감사드린다. 이론의 시험을 위한 귀중한 데이터를 제공한 V.E. Fortov와 V.V. Yakushev의 실험 그룹에 특별한 감사를 표한다.
	
	\section*{자금 지원}
	
	본 연구는 야쿠셰프 연구소의 지원을 받았다.
	
	\section*{데이터 가용성}
	
	모든 계산, 코드 및 보충 자료는 \url{https://github.com/Alexey-Yakushev-YUCT/YPSDC}에서 이용 가능하다.
	
	\begin{thebibliography}{99}
		
		\bibitem{fortov1999}
		V.E. Fortov, V.V. Yakushev, K.L. Kagan et al., "Anomalous electric conductivity of lithium under quasi-isentropic compression to 60 GPa (0.6 Mbar). Transition into a molecular phase?", JETP Letters, 70, 628 (1999).
		
		\bibitem{fortov2001}
		V.E. Fortov, V.V. Yakushev, K.L. Kagan, I.V. Lomonosov, V.I. Postnov, T.I. Yakusheva, A.N. Kuryanchik, "Anomalous resistivity of lithium at high dynamic pressure", Pis'ma v Zh. Eksper. Teoret. Fiz., 74:8, 458 (2001) [JETP Letters, 74:8, 418 (2001)].
		
		\bibitem{yakushev2002}
		V.E. Fortov, V.V. Yakushev et al., "Lithium at high dynamic pressure", Journal of Physics: Condensed Matter, 14, 10809 (2002).
		
		\bibitem{orlov2001}
		A.I. Orlov, L.G. Khvostantsev, E.L. Gromnitskaya, O.V. Stal'gorova, "Effect of pressure on kinetic properties and phase transitions in lithium", JETP, 120:2, 445 (2001).
		
		\bibitem{tsallis1988}
		C. Tsallis, "Possible generalization of Boltzmann-Gibbs statistics", Journal of Statistical Physics, 52, 479 (1988).
		
		\bibitem{bertulani2015}
		C.A. Bertulani et al., "Big Bang nucleosynthesis with a non-Maxwellian distribution", Journal of Physics G: Nuclear and Particle Physics, 42, 125201 (2015).
		
		\bibitem{west1997}
		G.B. West, J.H. Brown, B.J. Enquist, "A general model for the origin of allometric scaling laws in biology", Science, 276, 122 (1997).
		
		\bibitem{yuct2026}
		A.V. Yakushev, "Yakushev Unified Coordination Theory (YUCT)", Zenodo, \url{https://doi.org/10.5281/zenodo.18444599} (2026).
		
		\bibitem{appendixL}
		A.V. Yakushev, "Appendix L: Fractal Coordination Error Scaling – A Universal Law from DNA to Cosmology", YUCT (2026).
		
		\bibitem{appendixC}
		A.V. Yakushev, "Appendix C: The Coordination-Based Periodic Table of Elements and Experimental Verification of the Universal Scaling Law", YUCT (2026).
		
		\bibitem{appendixP}
		A.V. Yakushev, "Appendix P: Temperature Scaling of Coordination Efficiency in Molecular Systems", YUCT (2026).
		
		\bibitem{Crumpler2024}
		N.R. Crumpler et al., "Gravitational redshift of white dwarfs in SDSS DR1-19", \emph{Astrophysical Journal} 977, 237 (2024).
		
		\bibitem{Fontaine2001}
		G. Fontaine, P. Brassard, P. Bergeron, "The potential of white dwarf cosmochronology", \emph{Publications of the Astronomical Society of the Pacific} 113, 409 (2001).
		
		\bibitem{Lantink2022}
		M.L. Lantink et al., "Earth's longest fossil rift-valley system", \emph{Nature Geoscience} 15, 571 (2022).
		
		\bibitem{Farhat2022}
		M. Farhat et al., "The resonance capture of the Moon and the origin of its orbit", \emph{Astronomy \& Astrophysics} 665, A108 (2022).
		
		\bibitem{IAEA2017}
		F. Hammache et al., "New determination of the \(^3\mathrm{He}(\alpha,\gamma)^7\mathrm{Be}\) cross section and its impact on the cosmological lithium problem", \emph{EPJ Web of Conferences} 165, 01020 (2017).
		
		\bibitem{HAL2006}
		V.F. Dmitriev, V.V. Flambaum, "The cosmological lithium problem and the variation of fundamental constants", \emph{Phys. Rev. D} 74, 063514 (2006).
		
		\bibitem{Knizhnik1988}
		V.G. Knizhnik, A.M. Polyakov, A.B. Zamolodchikov, "Fractal structure of 2D quantum gravity", Mod. Phys. Lett. A 3, 819 (1988).
		
		\bibitem{Tsallis2009}
		C. Tsallis, "Introduction to Nonextensive Statistical Mechanics", Springer (2009).
		
	\end{thebibliography}
	
\end{document}